\documentclass[11pt,a4paper]{article}
\usepackage[T1]{fontenc}
\usepackage[utf8]{inputenc}
\usepackage[main=italian,provide=*]{babel}
\usepackage{amsmath,amssymb}
\usepackage{geometry}
\geometry{margin=2.4cm,top=2.6cm,bottom=2.8cm}
\usepackage{booktabs}
\usepackage[table]{xcolor}
\usepackage{tcolorbox}
\tcbuselibrary{skins,breakable}
\usepackage{titlesec}
\usepackage{enumitem}
\usepackage{needspace}
\usepackage{fancyhdr}
\usepackage{hyperref}
\hypersetup{colorlinks=true,linkcolor=Sepia,citecolor=Sepia,urlcolor=Sepia}

% --- palette ---
\definecolor{inkblue}{HTML}{1B3A5C}
\definecolor{lightblue}{HTML}{EAF1F8}
\definecolor{accent}{HTML}{B0562A}
\definecolor{softgray}{HTML}{6B6B6B}
\definecolor{okgreen}{HTML}{2E6B3E}
\definecolor{okgreenbg}{HTML}{EAF5EC}

% --- titoli di sezione ---
\titleformat{\section}
  {\normalfont\Large\bfseries\color{inkblue}}
  {\thesection}{0.6em}{}[{\color{inkblue!40}\titlerule}]
\titlespacing{\section}{0pt}{0.7em}{0.4em}
\titleformat{\subsection}
  {\normalfont\large\bfseries\color{accent}}
  {}{0em}{}
\titlespacing{\subsection}{0pt}{0.5em}{0.2em}

\pagestyle{fancy}
\fancyhf{}
\renewcommand{\headrulewidth}{0.4pt}
\fancyhead[L]{\small\color{softgray}\texttt{dimer\_exact.py} --- spiegazione e risultati}
\fancyhead[R]{\small\color{softgray}\thepage}
\fancyfoot[C]{\small\color{softgray}Tesi triennale, Samuele Schiavi --- Relatore: Prof.\ Alessandro Chiesa}

% --- box riutilizzabili ---
\newtcolorbox{keyeq}[1][]{
  colback=lightblue, colframe=inkblue!70, boxrule=0.6pt,
  arc=2pt, left=8pt, right=8pt, top=2pt, bottom=2pt,
  #1
}
\newtcolorbox{testbox}{
  colback=okgreenbg, colframe=okgreen!60, boxrule=0.6pt,
  arc=2pt, left=8pt, right=8pt, top=2pt, bottom=2pt,
  fonttitle=\bfseries\color{okgreen}, title=Output del self-test
}
\newtcolorbox{panelbox}[1]{
  colback=white, colframe=softgray!50, boxrule=0.5pt,
  arc=2pt, left=7pt, right=7pt, top=2pt, bottom=2pt,
  fonttitle=\bfseries\color{inkblue}, title=#1
}

\title{\vspace{-1.5em}\color{inkblue}\texttt{dimer\_exact.py}\\[0.2em]
\Large\normalfont\color{softgray}Spiegazione concettuale e lettura dei risultati}
\author{Note di lavoro --- Tesi triennale, Samuele Schiavi\\ Relatore: Prof.\ Alessandro Chiesa}
\date{}

\begin{document}
\sloppy
\maketitle
\thispagestyle{fancy}
\vspace{-2.6em}

\section*{Cosa fa il modulo}
\addcontentsline{toc}{section}{Cosa fa il modulo}

Calcola la soluzione esatta del dimero di spin-1/2 ($N=2$) per una
griglia di valori del campo $b$, restituendo spettro, ground state,
magnetizzazione e un self-test automatico contro la formula analitica.
Il notebook di supporto (\texttt{dimer\_exact\_spiegato.ipynb}) contiene
due modi distinti e complementari di ottenere lo spettro esatto:

\begin{panelbox}{Spettro analitico --- \texttt{analytic\_eigenvalues}}
La formula chiusa
\begin{equation}
E(S,M) = 2J\,S(S+1) - 3J + 2bM
\end{equation}
valida \textbf{solo per $D=0$}. È la soluzione algebrica vera, ottenuta a
carta e penna sfruttando le simmetrie ($[H_1,H_2]=0$, buoni numeri
quantici $S$ e $M$). È implementata in tre righe di Python puro, senza
diagonalizzare nulla.
\end{panelbox}
\smallskip

\begin{panelbox}{Diagonalizzazione numerica esatta --- \texttt{exact\_sweep} con \texttt{eigh}}
Costruisce la matrice $4\times4$ di $H$ e la diagonalizza numericamente.
È "esatta" nel senso che non usa nessuna approssimazione variazionale
(non è un VQE), ma è comunque un calcolo numerico --- soggetto a errori
di arrotondamento floating point (i $\sim10^{-16}$ del self-test qui
sotto). \textbf{Funziona anche per $D\neq0$}, dove la soluzione algebrica
chiusa non esiste.
\end{panelbox}

Il notebook contiene entrambe, e il self-test le confronta proprio per
verificare che coincidano (a meno di $\sim10^{-16}$). La diagonalizzazione
numerica è il \emph{benchmark} contro cui si misura il VQE nel resto
della tesi --- il riferimento "vero" rispetto al quale l'errore
variazionale è giudicato:

\begin{keyeq}
\centering
\footnotesize
\begin{equation*}
\underbrace{E(S,M) = 2JS(S+1)-3J+2bM}_{\text{algebra esatta }(D=0)}
\;\approx\;
\underbrace{\texttt{eigh}}_{\text{numerico esatto (}D\text{ qualunque)}}
\;\geq\;
\underbrace{\langle H\rangle_{\boldsymbol\theta}}_{\text{VQE (variazionale, approssimato)}}
\end{equation*}
\end{keyeq}

Il VQE cerca il minimo di $\langle H\rangle$ su una famiglia di stati
parametrizzati, e per il teorema variazionale
$\langle H\rangle_{\boldsymbol\theta} \geq E_0$ sempre --- può solo
avvicinarsi all'esatto, mai superarlo.

\section*{Verifica numerica: il self-test}
\addcontentsline{toc}{section}{Verifica numerica}

\begin{testbox}
\texttt{max\textbar{}E\_num - E\_analitico\textbar{} (D=0) = 8.88e-16}
\end{testbox}

$8.88 \times 10^{-16}$ è rumore di arrotondamento a precisione macchina
($\varepsilon_{\text{float64}} \approx 2.2 \times 10^{-16}$), otto ordini
di grandezza sotto la soglia di $10^{-10}$. Lo spettro numerico
\textbf{coincide} con quello analitico $E(S,M) = 2JS(S+1) - 3J + 2bM$:
nessuna discrepanza fisica.

\section*{Fin dove arriva la formula chiusa}
\addcontentsline{toc}{section}{Fin dove arriva la formula chiusa}

Si può confrontare il VQE contro la formula chiusa, ma solo per $D=0$ e
solo per il dimero $N=2$. La formula chiusa esiste perché il dimero
isotropo ha due simmetrie esatte che rendono $H$ diagonale nella base
$|S,M\rangle$: $[H,S^2]=0$ e $[H,S_z]=0$. È un caso eccezionalmente
fortunato. Non appena si esce da questo caso, la formula chiusa non
esiste più:

\begin{itemize}[leftmargin=1.4em, itemsep=0.15em]
\item \textbf{Con $D\neq0$}: $[H,S_z]\neq0$, i buoni numeri quantici si
rompono, lo spettro non si scrive in forma chiusa. Occorre diagonalizzare
numericamente.
\item \textbf{Con $N=3$} (catena o triangolo): l'Hamiltoniana è
$8\times8$; ci sono ancora simmetrie (es. $S_z$ totale per la catena), ma
lo spettro completo non ha una forma chiusa semplice.
\item \textbf{Con rumore (Parte 2)}: lo stato è $\rho$, non $|\psi\rangle$,
e il confronto va fatto sull'energia dell'operatore densità --- ancora
numerico.
\end{itemize}

\medskip
\noindent La struttura del confronto nella tesi è quindi:

\begin{center}
\renewcommand{\arraystretch}{1.3}
\begin{tabular}{@{}p{5.4cm}p{6.4cm}@{}}
\toprule
\rowcolor{inkblue!10}
\textbf{Sistema} & \textbf{Riferimento per il VQE}\\
\midrule
$N=2$, $D=0$ & formula chiusa \textbf{e} \texttt{eigh} (coincidono)\\
$N=2$, $D=0.2$ & solo \texttt{eigh}\\
$N=3$ (entrambe le topologie) & solo \texttt{eigh}\\
Parte 2 (rumore) & solo \texttt{eigh} (su $\rho$)\\
\bottomrule
\end{tabular}
\end{center}

Per il dimero con $D=0$ si usano entrambi --- ed è utile farlo
esplicitamente, perché mostra la catena di validazione completa:

\begin{keyeq}
\centering
formula chiusa $\approx$ \texttt{eigh} $\approx$ VQE \quad
(con errore variazionale misurabile)
\end{keyeq}

Il confronto operativo del VQE è però sempre contro \texttt{eigh},
perché è quello che scala a tutti i casi in roadmap.

\section*{Lettura della figura (\texttt{dimer\_exact.png})}
\addcontentsline{toc}{section}{Lettura della figura}

\begin{panelbox}{Pannello sinistro --- Spettro $E/J$}
\begin{itemize}[leftmargin=1.2em, itemsep=0.1em]
\item con $D=0$ (nero) il fondamentale è il singoletto ($E=-3J$, piatto)
fino a $B/J=2$, poi lo stato $M=-1$ che scende linearmente; l'incrocio è
netto;
\item con $D=0.2$ (rosso tratteggiato) il fondamentale passa con
continuità attraverso $B/J=2$: l'incrocio è diventato \textbf{evitato}
(anticrossing), con gap minimo $\Delta \simeq 0.57$.
\end{itemize}
\end{panelbox}
\smallskip

\begin{panelbox}{Pannello destro --- Magnetizzazione $\langle M_z\rangle$}
\begin{itemize}[leftmargin=1.2em, itemsep=0.1em]
\item con $D=0$: salto netto da $0$ (singoletto) a $-1$ (tripletto
$M=-1$) in $B/J=2$ --- firma della transizione del primo ordine;
\item con $D=0.2$: crossover liscio attraverso $B/J=2$ --- conseguenza
diretta dell'anticrossing.
\end{itemize}
\end{panelbox}

\section*{Ruolo nel progetto}
\addcontentsline{toc}{section}{Ruolo nel progetto}

\texttt{dimer\_exact.py} è il \textbf{metro di misura} dell'intera tesi. I
moduli VQE importano \texttt{dimer\_hamiltonian}, \texttt{exact\_sweep} e
\texttt{analytic\_eigenvalues} da questo file e non lo modificano mai. La
catena di validazione è, ancora una volta:

\begin{keyeq}
\centering
\begin{equation*}
\underbrace{E(S,M)}_{\text{formula chiusa}} \approx
\underbrace{\texttt{eigh}}_{\text{numerico esatto}} \geq
\underbrace{\langle H\rangle_{\boldsymbol\theta}}_{\text{VQE}}
\end{equation*}
\end{keyeq}

\end{document}
